\documentclass{article}
\usepackage{amsmath,amssymb,amsthm}
\newtheorem{proposition}{Proposition}
\newtheorem{conjecture}[proposition]{Conjecture}
\theoremstyle{definition}
\newtheorem{definition}[proposition]{Definition}
\newtheorem{example}[proposition]{Example}
\newtheorem{remark}[proposition]{Remark}
\newcommand{\taylor}[1]{\textbf{[Owner: #1]}}
\newcommand{\todo}[1]{\textbf{[To do: #1]}}
\newcommand{\agentreply}[1]{\textit{[Agent: #1]}}
\begin{document}
\section{Bounds}
% Human note: do not lose the interval restriction.
\taylor{Keep a worked value between the endpoints.}
% >>> [agent-fixture 2026-09-18 bounded-square]
The interval $[0,1]$ consists of real numbers $x$ with $0\leq x\leq1$.
\begin{proposition}\label{prop:bounded-square}
Let $x$ be a real number with $0\leq x\leq1$. Then $x^2\leq x$.
\end{proposition}
\begin{proof}
Both $x$ and $1-x$ are nonnegative. Thus $x-x^2=x(1-x)\geq0$.
\end{proof}
\begin{remark}\label{rem:bounded-square-equality}
For real $x$ in $[0,1]$, equality holds exactly at $x=0$ or $x=1$,
by the zero-product property applied to $x(1-x)$.
\end{remark}
For example, $x=1/2$ gives $x^2=1/4\leq1/2$.
% <<< [agent-fixture 2026-09-18 bounded-square]

\section{Closing note}
This paragraph belongs to the owner and is outside the current research.
\end{document}
